\section{Introduction}
\label{sec:introduction}

Bostrom's simulation argument~\cite{Bostrom2003are} made the idea that we might live in a computer simulation a topic for academic philosophy, and it has drawn a large critical literature (Sec.~\ref{sec:philosophy}). Physicists have since proposed ways to test the idea. They have looked for lattice artefacts in the highest-energy cosmic rays~\cite{Beane2014constraints}, proposed interferometric searches for Planckian ``pixelation'' of space~\cite{Hogan2008measurement,Chou2016first}, and suggested quantum-optics experiments to detect a simulator that renders only what is observed~\cite{Campbell2017on}. Others have argued that undecidability rules simulation out~\cite{Faizal2025consequences}, a complexity result on quantum Monte Carlo~\cite{Ringel2017quantized} that does not mention simulation at all was widely reported as doing so (Sec.~\ref{sec:methodology:reception}), and it has been argued that an information-theoretic ``second law of infodynamics'' supports it~\cite{Vopson2023second}. These results have often been reported as proving or disproving the hypothesis, beyond what the papers themselves claim (Sec.~\ref{sec:methodology:reception}).

The literature is scattered across philosophy, lattice field theory, astroparticle physics, quantum optics, complexity theory and thermodynamics, and we are not aware of a review that treats all of it together. This review has three aims:
\begin{enumerate}
\item to set out, for each proposal, what it assumes, what it derives, and what has actually been observed, with bounds, detections and proposals kept distinct;
\item to give a methodological framework that states precisely what any such test can and cannot establish (Sec.~\ref{sec:methodology});
\item to identify where new work is possible, and to carry out the cheapest parts of it.
\end{enumerate}

\paragraph{Position.} We do not argue for or against the simulation hypothesis. Our thesis is methodological. The unrestricted hypothesis is not an empirical claim. Only its conjunctions with explicit assumptions about a simulator's architecture, resources and non-adaptivity can be tested. Each such conjunction predicts signatures that non-simulation physics also predicts (Table~\ref{tab:subhyp}). Throughout, we therefore speak of \emph{constraining simulator designs}, never of proving or disproving the simulation hypothesis.

\paragraph{Contributions of this work.} Besides the review itself, the paper contains analyses that are, to our knowledge, new. Each is labelled ``this work'' where it appears.
\begin{itemize}
\item Mapping the lattice dispersion relations of Beane et al.~\cite{Beane2014constraints} onto published Lorentz-invariance-violation bounds. For their unimproved-Wilson scenario this gives $b^{-1}\gtrsim10^{16}$~GeV from the Crab synchrotron spectrum~\cite{Liberati2012scale}, compared with the original $\sim10^{11}$~GeV. We also show that $\gamma\to e^+e^-$ is kinematically allowed far below the lattice scale, and we determine how the bounds change under improved discretizations (Sec.~\ref{sec:lattice-liv:mapping}).
\item A documented search showing that no analysis of cosmic-ray arrival directions for a cubic-symmetric anisotropy has been published, and a specification of such an analysis using public Pierre Auger Observatory data (Sec.~\ref{sec:uhecr:proposal}).
\item An analysis showing that the success criteria of the proposed render-on-demand experiments~\cite{Campbell2017on} contradict standard quantum mechanics by up to $1/2$ in probability (Sec.~\ref{sec:rendering-qm-campbell}).
\item A reproduction of the SARS-CoV-2 entropy values underlying the ``second law of infodynamics''~\cite{Vopson2022second}. It shows that the reported decrease is accounted for by the virus's known mutational bias. We also give a calorimetric consistency test of the mass--energy--information conjecture (Sec.~\ref{sec:information}).
\end{itemize}

\paragraph{Verification protocol.} Referees and readers should be able to check every statement. Each factual claim about the literature is recorded in a claim ledger. The ledger gives the source, a precise locator (section, equation, page or table) and a verbatim supporting quotation, and it records whether we read the full text, only the abstract, or only the bibliographic record. Claims we could not verify are excluded from the text. Bibliographic data were taken from DOI, INSPIRE-HEP or arXiv metadata. A script checks that every citation is backed by a verified ledger entry and that every DOI and arXiv identifier resolves to the cited title. The ledgers, derivations and reproduction scripts are released with the paper (Sec.~\ref{sec:discussion:limitations}).

\paragraph{Outline.} Section~\ref{sec:philosophy} reviews the simulation argument and its critics, and Sec.~\ref{sec:digital-physics} separates the simulation hypothesis from neighbouring ``digital physics'' theses. Section~\ref{sec:methodology} gives the methodological framework. Sections~\ref{sec:lattice-liv}--\ref{sec:information} review the proposed signatures: lattice discreteness and Lorentz violation, ultra-high-energy cosmic rays, holographic noise, complexity and computability, render-on-demand quantum measurement, and information physics. Section~\ref{sec:discussion} summarizes what is constrained and sets out open problems.
